% !TeX root = ../papers/nc-p3.tex
\section{反対八次式の重複平面と三重点を持つ六次残余}
\label{sec:repeated-octic-plane-residual}
\begin{lemma}\label{lem:quintic-six-marked-triples}
四点が共線でない被約な印のうち、整五次曲面の三重以上の点は高々六つである。
六点なら二本の素な直線に三点ずつ載る。両直線を入力の二座標組 $u,v$ の軸とすると、五次式は
\[
                   P_3(u)Q_2(v)+P_2(u)Q_3(v)
\]
であり、添字は次数、三次式は相異なる三根を持つ。両軸の外の特異点は存在しても一本の直線に入る。
\end{lemma}
\begin{proof}
二つの三重点を結ぶ直線は五次曲面に含まれる。一平面内に一般位置の四点があれば
六本の割線が平面五次切断に入り不可能である。平面内に共線な三点があればその直線は少なくとも二重である。
一度除いても残余四次に三つの二重点があるためである。この直線外にさらに二点あれば、
外の各点から三点への計六本のうち一致は多くても一対であり、やはり次数を超える。
従って各平面では高々四点、特に二本の三点直線は交わることができない。

非共面な四点を四座標点に取り、標準三次Cremona変換を用いる。
座標点での重複度を $m_i\ge3$ とすると、変換後の整な真変換の次数は
$15-\sum m_i\le3$ である。実際五次の各単項式に対して四座標をそれぞれ他の三座標の積で置換すると、
第 $i$ 座標の最小指数が $m_i$ となる。座標平面外は同型であり、残りの三重点はその開集合では高々一つである。
整三次が異なる二点で三重なら両点を頂点とする錐となって二変数三次式に分解してしまう。
座標面内の追加点はその面の三頂点のどれか二つを結ぶ辺上に限る。各辺に一つまでで、
三点直線同士が素であるから使える辺は高々二本。合計七点以下である。
七点なら二本の三点直線とそれらの外の一点がある。第三の三点直線にはさらに三点必要で存在しない。
外の一点、第一直線の二点、これらの平面外の第二直線の一点を新たな座標点とすると、
追加点を持てる辺は第一直線だけになり七点を収容できない。
六点なら任意の四点基底の少なくとも一辺に三点が必要である。三点直線が一本だけなら
残る三点は非共線で、それらの平面外の直線上の一点を加えて基底にできるが、三点を持つ辺がない。
従って六点の配置は主張の通りである。

二軸に関する二次数 $(5-j,j)$ に分ける。一方の軸の三点で三重になるには $j=0,1$ の
係数二項式はそれぞれ三点で少なくとも三重、二重であり恒等零である。
$j=2$ の係数は三つの指定根を持つ二項三次の倍である。逆側も同様で表示を得る。
整性より $P_3,P_2$ は互いに素で $Q_3,Q_2$ も互いに素。
両軸の外の特異点では二群のEuler式により両項が別々に零である。
$P_3=Q_3=0$ では両根が単純なので滑らかである。
$P_2=Q_2=0$ も一方が単純なら滑らかで、両方二重のときだけ対応する一対の方向の横断直線が残る。
\end{proof}

\begin{lemma}\label{lem:normal-surface-nonrational-point-count}
標数零の整な正規射影超曲面 $S\subset\mathbb P^3$ が次数 $e\ge3$ なら、
非有理特異点の個数は $1+\binom{e-1}{3}$ 以下である。
\end{lemma}
\begin{proof}
個数を $b$、表示の上界を $c=\chi(\mathcal O_S)$ とし、$b>c$ とする。
滑らかな解消 $f:Y\to S$ では $\chi(\mathcal O_Y)=c-\sum_p p_g(p)\le c-b<0$。
従って滑らかな射影曲面の分類により $Y$ は種数 $g\ge2$ の曲線上の線織曲面に双有理である。
必要ならさらに解消して射影直線束の点吹上げとしての射 $Y\to B$ を取る。
非有理特異点の例外集合は必ず $B$ を支配する曲線を含む。実際全てが垂直なら連結性により一繊維 $F$ に入り、
各形式近傍の例外スキームはある $mF$ に含まれる（台の包含とNoether性）。
射影直線束の点吹上げでは $R^1\pi_*\mathcal O_Y=0$、$\pi_*\mathcal O_Y=\mathcal O_B$ で、
$0\to\mathcal O_Y(-mF)\to\mathcal O_Y\to\mathcal O_{mF}\to0$ を押し出すと
$H^1(\mathcal O_{mF})=0$ を得る。一次元スキームの商にも $H^1$ の消滅が降り、
形式関数定理により $R^1 f_*\mathcal O_Y$ がその点で零となってしまう。
水平な例外曲線の正規化の種数はRiemann--Hurwitzにより $g$ 以上である。
その $H^1$ は曲線の構造層の $H^1$ の商であり、各無限小例外近傍の $H^1$ からも全射を受ける。
これらの遷移写像は全射なので形式関数定理から $p_g(p)\ge g$ を得る。
従って $bg\le\sum p_g(p)=c+g-1$、すなわち $(b-1)g\le c-1$ となり $b>c$ に反する。
\end{proof}

\begin{theorem}\label{thm:repeated-octic-plane-residual}
被約な四辺有限窓の一般可逆な全階数反対根について、行列式が重複因子を持つなら
$h^2G$ であり、$h$ は平面、$G$ は $h$ とは素な被約六次式である。
$G$ は非正規な整六次式か、整四次因子と次数二の残余の積に限られる。
\end{theorem}
\begin{proof}
重複因子は平面 $H=(h=0)$ 一つだけである。$h$ を最大限除いた残余は
$H$ 外の少なくとも十三印の各点で三重以上である。被約四次式の三重印は五つ以下である。
整四次では三重点は共線、$(3,1)$ では整三次の特異印四つ以下とその平面外の三重頂点一つまで、
$(2,2)$ では二頂点までである。$(2,1,1)$ では二平面の交線と二次の頂点に限られ、
四平面では四つの三重交点か、三平面以上が含む一本の直線に限られる。
低い次数の被約式にも一般平面を加えてこの上界を適用できる。
被約五次式に平面があればその七印と残余四次の五印で計十二以下、
$(3,2)$ なら両因子の特異印で五つ以下、整なら前補題で六つ以下である。
従って $h$ の指数はちょうど二である。

$G$ が整かつ正規なら十三以上の三重印は非有理特異点である
（Gorenstein有理曲面特異点は有理二重点）。前補題の六次上界十一に反する。
五次因子と平面 $P$ の分割なら $H\cup P$ の外の少なくとも六印が五次の三重印となる。
六印で丁度であり、$H,P$ は相異なる七印ずつを含み、両方に属する印はない。
前補題の二本の軸はこの六印で直線の上限に達する。$P$ 上の七印は全て五次の二重以上の点で、
軸の外にある。そこでは一本の直線にしか特異点を持てず矛盾する。
二つの整三次の積では各特異集合に四印まで、合計八以下となる。
$(3,2,1)$ なら平面の七印と、その外で三次または二次因子の特異点にある五印で十二以下である。
$(3,1,1,1)$ では $h^2G$ の全二十印での重複度和は平面一枚につき七、
整三次につき二十四の上界から $5\cdot7+24=59$ 以下で、必要な六十に反する。
次数二以下の因子だけでも整二次の重複度上界十五で同様に $2\cdot7+3\cdot15=59$ 以下である。
以上で全ての分割を確認した。
\end{proof}

\begin{corollary}\label{cor:repeated-octic-plane-integral-sextic}
実際、重複平面を二回除いた残余は常に整かつ非正規な六次式である。
\end{corollary}
\begin{proof}
$G=FQ$ とし、$F$ は整四次、$Q$ は被約二次とする。三次式 $hQ$ は非零なので三次評価の単射性により
その外に一つの印 $p$ がある。この点は $F$ の三重以上の点である。
このような点が一つあれば $F$ の特異印は全部で七以下であることを示す。
$p$ を座標点にし、四次式を $t A_3(z)+A_4(z)$ とする。
$p$ が四重点なら $A_3=0$ で整平面四次曲線上の錐である。その特異方向は種数公式により三つ以下で、
各母線には $p$ の外に二印までなので計七以下である。
それ以外なら $A_3,A_4$ は互いに素で、$p$ の外の特異点の方向はその交点である。
実際 $t$ 微分と四次式から両者が消え、特異点ではさらに $t\,dA_3+dA_4=0$ である。
方向に特異点が一つあれば二つの平面曲線の局所交数は二以上である。
同じ方向に二つの異なる特異点があれば $dA_3=dA_4=0$ となり交数は四以上で、
この母線には $p$ 以外に二印までである。
よって全交数十二に対して $p$ 外の特異印一つごとに少なくとも二を使い、主張を得る。
$Q$ の特異印は頂点一つか二平面の交線上の三つ以下である。
$G$ の三重点では少なくとも一方の因子が特異であるから、平面外の十三印以上に対し
収容できるのは $7+3=10$ 印までとなる。
\end{proof}
